\documentclass[11pt]{article}
\usepackage[a4paper,margin=1in]{geometry}
\usepackage{amsmath,amssymb,amsthm}
\usepackage[hidelinks]{hyperref}
\newtheorem{proposition}{Proposition}
\newtheorem{criterion}{Conditional criterion}
\title{REAL 1B: Gold 4.40-J/X\\Signed Ledger and Conditional Twin-Prime Interfaces}
\author{Provisional AI-assisted research addendum}
\date{21 November 2026}
\begin{document}
\maketitle
\begin{abstract}
This is a research-status supplement, not a proof of the Twin Prime Conjecture.
We document two elementary identities, a pointwise rough-prime minorant,
and the new source-specific arithmetic estimates still required to obtain a
strictly positive twin-prime reserve. Imported asymptotic assertions in the
Gold 4.39-A/4.40-J/4.40-X reports are subject to independent verification.
\end{abstract}
\section{Exact signed source identity}
Let \(L=\log X\), \(N_0=\mathfrak S_2X/L\), \(z=(2X)^{1/4}\) and
\[
a(n)=\mathbf1_{X<n\le2X}\vartheta_{\mathbb P}(n-2),\qquad
\vartheta_{\mathbb P}(m)=\log m\,\mathbf1_{m\in\mathbb P}.
\]
Write \(R_z(n)=\mathbf1_{P^-(n)>z}\), and assume
\[
\ell_B(n)=\alpha_XR_z(n)-w_X(n),\quad
0<\alpha_X\le1,\quad w_X\ge0,\quad w_XR_z=0.
\]
Let \(T,C_2,C_3\) denote the masses of squarefree rough \(n\) with
respectively one, two, or three prime factors, each weighted by \(a(n)\).
There are no higher squarefree rough orders on this interval.
Set \(O=(\mu^2-\mu)/2\), \(E=(\mu^2+\mu)/2\), and
\[
W_o=\sum a(n)w_X(n)O(n),\quad W_e=\sum a(n)w_X(n)E(n).
\]
For any auxiliary mark \(\mathcal W\) vanishing on rough inputs put
\[
A_{\rm sf}=\sum a\mu^2\ell_B,\quad
\delta=\sum a\mu\ell_B\mathcal W,\quad
\mathcal E_{\mathcal W}=\sum a\mu\ell_B(1-\mathcal W).
\]
\begin{proposition}[Exact ledger]
\begin{align*}
A_{\rm sf}&=\alpha_X(T+C_2+C_3)-W_o-W_e,\\
\mathcal E_{\mathcal W}+\delta
 &=\alpha_X(-T+C_2-C_3)+W_o-W_e,\\
T&=\frac{A_{\rm sf}-\mathcal E_{\mathcal W}-\delta+2W_o}
 {2\alpha_X}-C_3.
\end{align*}
\end{proposition}
\begin{proof}
On rough squarefree inputs \(\mu=-1,+1,-1\) according to the
one-, two-, and three-prime-factor cases. On nonrough inputs
\(\ell_B=-w_X\), yielding \(-W_o-W_e\) in the unsigned sum and
\(W_o-W_e\) in the signed sum. Subtract the resulting equalities.
\end{proof}
Define \(c_B=\frac72\log3-4\log2\), \(c_3=0.49073558986414\ldots\)
and \(h_0=c_B-2c_3=0.0910831083703276\ldots\).
\begin{criterion}[Unmarked H2: OPEN]
Assume the inherited source inequalities
\(A_{\rm sf}\ge(c_B-e-o(1))N_0\),
\(C_3\le(c_3+o(1))N_0\), \(\delta=o(N_0)\),
and \(\alpha_X\ge\alpha_0>0\).
If the as-yet \emph{unproved} estimate
\[
\mathcal E_{\mathcal W}\le(h_0-e-\kappa)N_0
\quad(\kappa>0\text{ fixed})
\]
also holds, then
\(
T/N_0\ge c_3(\alpha_X^{-1}-1)+\kappa/(2\alpha_X)-o(1)>0.
\)
\end{criterion}
\begin{proof}
Insert the assumptions in the exact ledger and use \(W_o\ge0\).
\end{proof}
Gold 4.39-A reports a marked Type-II estimate
\(\mathscr V_{u,e}[\mathcal W]=O_A(XL^{-A})\) with imported operator
hypotheses, but \(\mathcal W\) vanishes on the rough sector.
It does \emph{not} prove the unmarked estimate above.
\section{Gold 4.40-X pointwise minorant}
Fix \(3/8<\eta<1/2\) and put
\[
b_\eta(n)=\frac1L\sum_{\substack{p\mid n\\p>X^\eta}}\log p,
\quad g_\eta(n)=\mu(n)b_\eta(n),\quad
s=1-\eta,\quad t_X=\frac{\log(2X)}{2L},
\]
\[
\lambda_X=\frac{s-t_X}{2}>0,\qquad
a_X^*=\frac{s+t_X}{2}.
\]
\begin{proposition}[Pointwise rough prime minorant]
For every \(X<n\le2X\),
\[
R_z(n)\{\lambda_X\mu^2(n)+\mu(n)[a_X^*-b_\eta(n)]\}
\le t_X\mathbf1_{n\in\mathbb P}.
\]
\end{proposition}
\begin{proof}
For a rough prime the left side is \(b_\eta-t_X\le t_X\).
For a rough semiprime at least one factor exceeds \(X^\eta\),
and \(b_\eta\ge1-\eta=s\); the left side is \(s-b_\eta\le0\).
For a rough triple at most one prime exceeds \(X^\eta\), since
\(2\eta+1/4>1\). If present, it is below \(\sqrt{2X}\),
so \(b_\eta\le t_X\); the left side is \(b_\eta-t_X\le0\).
On nonsquarefree inputs both \(\mu,\mu^2\) vanish.
Four rough prime factors would exceed \(2X\).
\end{proof}
Let
\[
S_X=\sum a(n)\mu^2(n)R_z(n),\quad
M_X^\mu=\sum a(n)\mu(n)R_z(n),\quad
G_X=\sum a(n)g_\eta(n)R_z(n).
\]
Summing the pointwise inequality against \(a(n)\ge0\) yields
\[
t_XT\ge\lambda_XS_X+a_X^*M_X^\mu-G_X.
\]
Furthermore \(S_X\ge A_{\rm sf}\) follows from the signed kernel
decomposition with \(0<\alpha_X\le1\) and \(w_X\ge0\).
\begin{criterion}[Paired-prime signed input: OPEN]
Assume \(A_{\rm sf}\ge(c_B-e-o(1))N_0\).
If a new estimate on this \emph{same two-prime source} proves
\[
G_X-a_X^*M_X^\mu
 \le[\lambda_X(c_B-e)-\kappa]N_0\quad(\kappa>0),
\]
then \(T\ge(\kappa/t_X-o(1))N_0>0\) on unbounded scales.
\end{criterion}
\begin{proof}
Insert the proposed bound and the inherited lower bound in the
pointwise inequality; \(t_X\to1/2\).
\end{proof}
\section{Audit and scope}
Gold 4.40-J's absolute-Cauchy certificate satisfies
\(M=A_{\rm sf}+2W_0\), \(V=A_{\rm sf}+W_0+D_w\), and hence
\(\sqrt{MV}\ge A_{\rm sf}\). This is a restriction on \emph{that
upper-bound method}, not a lower bound on the actual signed residual.
The supplied J/X reports contain further claimed analytic estimates
which have not been independently re-established in this addendum.
Neither J's ordinary BV cofactor treatment nor X's ordinary-integer
distribution estimates imply the required \(\vartheta_{\mathbb P}(n-2)\)
signed contrast. Gold 4.39-B even minor arcs and the third Gold 4.40
source-capacity branch remain open at the last supplied checkpoint.
The two conditional criteria overlap and cannot be added as
independent arithmetic credit. No G3 historical novelty certification,
G5 twin-prime closure or Lean certification is claimed.
\begin{thebibliography}{9}
\bibitem{FM} K. Ford and J. Maynard, \emph{On the theory of prime producing sieves},
\url{https://arxiv.org/abs/2407.14368}.
\bibitem{LWW} G. L\"u, Z. Wang and J. Wu,
\emph{On the Hardy--Littlewood--Chowla conjecture},
\url{https://arxiv.org/abs/2610.20242}.
\bibitem{JLL} Y. Jiang, Y.-K. Lau and G. L\"u,
\emph{An asymptotic formula for primes of the form \(u^2+v^2+1\)},
\url{https://arxiv.org/abs/2611.12467}.
\end{thebibliography}
\end{document}